\section{Methodology}
\label{sec:method}



\paragraph{Workload.}
We use the public DAB release \cite{ma2026dab}: 13,482 of its 13,500 runs of five LLMs on 54 tasks over 12 datasets.
The LLMs are Gemini-2.5-Flash, Gemini-3-Pro, GPT-5-mini, GPT-5.2 and Kimi-K2, and the databases run on SQLite,
DuckDB, PostgreSQL and MongoDB. Parsing the trajectories yields 151,293 tool calls, of which 76,213 are
\texttt{query\_db} calls. Re-running each task's own validator on the final answers gives 4,543 passing and 8,939
failing runs. Table~\ref{tab:datasets} lists the datasets, which range from 0.6\,MB to 5.6\,GB on disk. The traces
record what each call returned but not what it cost, so we measure cost by replay.

% Table -- the 12 datasets (full width, Section 3)
\begin{table*}[tbp]
  \caption{The 12 DAB datasets of the workload. Columns give the engines (S: SQLite, D: DuckDB, P: PostgreSQL, M:
  MongoDB), the tasks, and the size on disk of the database files and dumps as shipped. They also give the
  \texttt{query\_db} calls, the distinct (database, query) pairs, and the share of successfully replayed calls whose
  result spills to a file. Isolated DB-execution time and payload bytes are summed over all calls. The last columns
  give the median over the dataset's tasks of the DB-time share repeated from other attempts and within the session
  (potential tier).}
  \label{tab:datasets}
  \small
  \begin{tabular}{@{}llrrrrrrrrr@{}}
\toprule
 & & & Data & & Distinct & Spilled & DB time & Bytes & \multicolumn{2}{c}{Task-median repeat \%} \\ \cmidrule(l){10-11}
Dataset & Engines & Tasks & (MB) & Calls & pairs & \% & (s) & (GB) & other attempts & session \\
\midrule
agnews & S, M & 4 & 40.6 & 4,240 & 1,360 & 32.8 & 196 & 11.07 & 41.8 & 2.3 \\
bookreview & S, P & 3 & 1.7 & 3,060 & 1,093 & 49.6 & 5 & 0.11 & 51.8 & 4.3 \\
crmarenapro & S, D, P & 13 & 32.8 & 16,879 & 10,412 & 11.7 & 73 & 0.39 & 35.5 & 1.8 \\
DEPS\_DEV\_V1 & S, D & 2 & 549.9 & 3,125 & 1,281 & 57.1 & 1,874 & 76.06 & 22.3 & 6.9 \\
GITHUB\_REPOS & S, D & 4 & 928.5 & 6,061 & 3,159 & 35.8 & 7,261 & 222.55 & 54.4 & 9.8 \\
googlelocal & S, P & 4 & 0.6 & 3,779 & 1,482 & 16.0 & 5 & 0.03 & 53.9 & 5.0 \\
music\_brainz\_20k & S, D & 3 & 3.1 & 2,451 & 1,071 & 23.2 & 32 & 1.67 & 50.0 & 0.8 \\
PANCANCER\_ATLAS & D, P & 3 & 301.4 & 6,716 & 4,045 & 43.4 & 211 & 1.05 & 37.9 & 3.6 \\
PATENTS & S, P & 3 & 5,556.3 & 4,263 & 2,322 & 63.0 & 23,398 & 225.08 & 36.1 & 7.1 \\
stockindex & S, D & 3 & 4.5 & 2,663 & 1,214 & 25.8 & 64 & 2.70 & 30.7 & 0.6 \\
stockmarket & S, D & 5 & 965.7 & 14,481 & 10,767 & 14.6 & 5,398 & 0.11 & 14.4 & 1.5 \\
yelp & D, M & 7 & 3.7 & 8,495 & 3,410 & 17.2 & 36 & 0.14 & 52.4 & 7.9 \\
\midrule
All &  & 54 & 8,388.9 & 76,213 & 41,616 & 25.9 & 38,553 & 540.96 & 41.2 & 3.2 \\
\bottomrule
\end{tabular}

\end{table*}

% Table -- workload, replay coverage and fidelity by engine (column width, Section 3)
\begin{table}[htbp]
  \caption{Replay coverage and fidelity by engine. Columns give the distinct (database, query) pairs and the share of calls
  whose replay succeeded. They also give the share of successful calls whose replayed result is byte-identical to
  the recorded observation. Byte mismatches differ only in row order or otherwise, and load-time object identifiers
  explain 1,299 of the 1,757 MongoDB mismatches. For spilled results, fidelity compares the 10,000-character preview that the trace keeps.
  The six pairs whose database is not in the task's configuration were not replayed.}
  \label{tab:workload}
  \small\setlength{\tabcolsep}{2.6pt}
  \begin{tabular}{@{}lrrrrrr@{}}
\toprule
 & & Distinct & Replay & Byte- & \multicolumn{2}{c}{Mismatches} \\ \cmidrule(l){6-7}
Engine & Calls & pairs & ok \% & identical \% & order & other \\
\midrule
DuckDB & 29,145 & 18,569 & 94.0 & 91.7 & 1,743 & 520 \\
MongoDB & 8,508 & 3,292 & 96.4 & 78.5 & 0 & 1,757 \\
PostgreSQL & 19,033 & 11,712 & 88.4 & 94.3 & 547 & 406 \\
SQLite & 19,455 & 8,037 & 98.5 & 99.9 & 0 & 4 \\
\midrule
All & 76,213 & 41,616 & 93.9 & 93.0 & 2,290 & 2,687 \\
\bottomrule
\end{tabular}

\end{table}


\paragraph{Isolated replay and fidelity.}
The 76,213 calls contain 41,616 distinct (database, query) pairs, and 66 calls lack a complete database or query
argument. We execute each pair three times in isolation in a version-pinned runtime through DAB's own execution and
serialization code. The runtime is a container image with Python 3.12, DuckDB 1.3.1 restricted to one thread and
pandas 2.3.0. It also holds PostgreSQL 16.15 without parallel workers and MongoDB 7.0.43, and it applies a 120\,s cap.
A pair is not executed again after an execution that reaches the cap. Six pairs name a database that is not in the
task's configuration and cannot be replayed. Hence 41,610 pairs covering 76,141 calls have a replay outcome
(Table~\ref{tab:workload}). Each pair yields an isolated execution time, which is the median of the completed runs.
It also yields serialization and spill times, the size of the serialized result, and a byte-stability flag over the
three results. In all, 37,675 pairs are stable, 3,917 fail, 15 time out and 3 return different values. We compare
each replayed result with the observation recorded in the original run. That observation is the full result when it
was inline and its 10,000-character preview when it was spilled. A replay is byte-identical to it, or equivalent
after normalization. Normalization covers row order, missing-value rendering, float digits beyond the ninth decimal
and MongoDB object identifiers assigned at load time. The two checks answer different questions. Stability over the
three replays is our empirical evidence that reuse is result-preserving inside the pinned runtime. The comparison
with the original run establishes \emph{fidelity}: the replayed cost and size describe the call the agent made. For a
spilled result it covers the preview, because the traces do not keep the file.

\paragraph{Cost model.}
We report two time metrics. \emph{Isolated DB-execution time} is the replayed execution time of a call and measures
the database work that reuse can avoid. \emph{Tool time} is what a design spends per call at the tool boundary. It
covers execution and serialization on a miss, the spill-file write, a cache lookup, and the per-file work on a hit.
The simulator charges the per-operation costs of the first live calibration. A lookup takes 1.27\,\textmu s, an alias
12.6\,\textmu s and a copy 0.46\,ns per byte. A clone, measured in the copy-on-write calibration, takes 0.29\,ms. Two
later calibrations measured the first three costs again and differ from these values by at most 7.3\%. When we
validate the simulator against a live run, it charges that run's costs. A simulation contains the calls that its tier
admits, and savings are relative to executing those calls without reuse. A task's saving is the ratio of its costs
summed over its models and resamples, and we report the median over tasks.

\paragraph{Admission tiers.}
A cache needs an admission rule that decides, before it reuses a result, that reuse is result-preserving. We evaluate
three tiers of decreasing strength (Table~\ref{tab:tiers}). \emph{Certified (Contract B)} is a default-deny whitelist
over the query text and the audited schema, a conditional certificate. Table~\ref{tab:contractb} summarizes its
construct whitelist, and the artifact holds the executable rule with its additional fail-closed checks. The rule
admits one \texttt{SELECT} block over base tables in which every name binds to a stored column. It admits three shapes. The first is one row of counts and extrema. The second is stored columns under an \texttt{ORDER BY} that names every output. The third is groups ordered on their grouping columns. \texttt{SUM}, \texttt{AVG} and the variance family are not on the list, because
their floating-point value depends on the order of accumulation
\cite{mueller2018reproducible,duckdb2026nondeterminism}. Clock-dependent values are excluded as well. On MongoDB,
whose DAB interface has no sort, the rule admits only a lookup by object identifier. The rule certifies byte
stability under three preconditions. P1 requires a fixed snapshot, engine version and session configuration, and a
fresh read-only session. P1 also requires a catalog without objects that act behind the query text, such as
row-level security or user-defined functions. P2 concerns every stored column that the query outputs, groups by or
passes to \texttt{MIN} or \texttt{MAX}. In such a column, values that the engine compares as equal print the same
bytes. P3 requires a deterministic serializer. Under P1--P3, two completing executions of an admitted query return
the same bytes whatever the plan, parallelism or physical row order. The joined and filtered rows form a multiset
that the text and the data determine, and \texttt{COUNT} is a function of that multiset. \texttt{MIN}, \texttt{MAX}
and the grouping range over columns for which P2 holds. The \texttt{ORDER BY} leaves ties only between rows that
print the same. The plan does not decide whether the query raises an error either. Every operation that
Table~\ref{tab:contractb} admits in a row expression returns a value or \texttt{NULL} for every input in the pinned
engines. A one-row output is evaluated exactly once. Contract B thus certifies the bytes of an execution that
completes. Whether a fresh execution completes within its resource limits depends on the query, the data, the limits
and the load. Completion therefore lies outside the certificate. P2 is a property of the data that the query text
cannot establish. A column that holds both $0.0$ and $-0.0$, or both $1$ and $1.0$, violates it. We therefore audit
the snapshot: all 19,772 stored columns of DAB's 26 SQL databases satisfy P2, and their catalogs hold no such object.
A unit suite of 181 assertions tests the implementation. Section~\ref{sec:threats} reports how the rule evolved after
the experiment plan and how its argument was checked.

\emph{Assumption-based (Contract E)} keeps the read-only, determinism and allow-list checks and drops the ordering
and shape checks. Byte stability then also rests on the runtime reproducing the engine's physical row order. Without a sort, the engines leave that order unspecified \cite{postgres2026order,sqlite2026select,mongodb2026sort}. We pin it through the engine version, one thread and a static snapshot. \emph{Potential (oracle)} is the population on
which we measure repeated work rather than an admission rule. It contains every call whose replay is equivalent to
its recorded observation, so that the replayed cost and size describe the call the agent made. It bounds what an
exact cache could reuse and cannot be computed before execution. Counting only byte-identical replays leaves the scope shares almost unchanged (Section~\ref{sec:robust}). Both static rules
are decided before execution. To measure designs, we restrict each rule to the calls whose replay succeeded, is
byte-identical to the recorded observation and is stable over three runs. These are the \emph{B-eval} and
\emph{E-eval subsets}. They are post-hoc evaluation restrictions that apply the same gates to both rules, so the two
subsets differ only in the rule.

% Table -- admission tiers (column width, Section 3)
\begin{table}[htbp]
  \caption{Admission tiers and the calls each admits (of 76,213). ``Online'' means decidable before the call executes.
  The B-eval and E-eval subsets keep the calls of each static rule whose replay was byte-identical to the recorded
  observation and stable over three replays. They are post-hoc restrictions on which we measure designs.}
  \label{tab:tiers}
  \small\setlength{\tabcolsep}{3pt}
  \begin{tabular}{@{}L{1.55cm}L{3.75cm}cr@{}}
\toprule
Tier & Admission rule & Online & Calls \\
\midrule
Certified & Contract B: conditional static byte-stability certificate & yes & 5,697 \\
 & \quad B-eval subset (post hoc) & no & 5,695 \\
Assumption-based & static Contract E, byte stability from a pinned runtime & yes & 74,894 \\
 & \quad E-eval subset (post hoc) & no & 66,263 \\
Potential & oracle: replay equivalent to the recorded observation & no & 70,173 \\
\bottomrule
\end{tabular}

\end{table}

% Table -- Contract B in full (column width, Section 3)
\begin{table}[htbp]
  \caption{Summary of Contract B's construct whitelist. Admission also applies the schema-binding, type-family,
  argument-form, alias and potentially-raising-expression checks of the artifact. $c$: a stored column that satisfies P2.}
  \label{tab:contractb}
  \footnotesize\setlength{\tabcolsep}{3pt}
  \begin{tabular}{@{}L{1.25cm}L{6.7cm}@{}}
  \toprule
  Statement & one \texttt{SELECT} block; no \texttt{WITH}, subquery, set operation, window, \texttt{DISTINCT ON},
  \texttt{VALUES} or table function \\
  \texttt{FROM} & base tables of the audited schema, joined by \texttt{JOIN}\,\ldots\,\texttt{ON} or a comma; no
  \texttt{USING}, \texttt{NATURAL} or \texttt{LATERAL} \\
  Names & each column binds to exactly one stored column of the tables in scope (ASCII names, the engine's case
  rule); an output alias does not shadow another column \\
  Row\newline expression & in \texttt{WHERE}, \texttt{ON} and the argument of \texttt{COUNT}: columns; literals without
  clock words; comparisons; \texttt{AND}, \texttt{OR}, \texttt{NOT}; \texttt{IS NULL}; \texttt{IN} over a list;
  \texttt{BETWEEN}; \texttt{LIKE} with a literal pattern; searched \texttt{CASE}; \texttt{COALESCE};
  \texttt{LOWER}, \texttt{UPPER}, \texttt{LENGTH}, \texttt{TRIM}, \texttt{REPLACE}. SQLite also: arithmetic,
  \texttt{ROUND}, \texttt{SUBSTR}, casts. DuckDB also: \texttt{SUBSTR}, \texttt{TRY\_CAST}, and arithmetic,
  \texttt{ABS} and \texttt{ROUND} over floating-point columns; a string meets only strings in a comparison \\
  One row & no \texttt{GROUP BY}, \texttt{HAVING}, \texttt{ORDER BY}, \texttt{OFFSET} or \texttt{DISTINCT}; each
  output built from \texttt{COUNT}, \texttt{MIN}($c$), \texttt{MAX}($c$), literals, arithmetic and the functions above \\
  Ordered & no aggregate; each output is a column $c$; \texttt{ORDER BY} names every output \\
  Grouped & \texttt{GROUP BY} columns $c$; each output is a grouping column or one \texttt{COUNT},
  \texttt{MIN}($c$) or \texttt{MAX}($c$); \texttt{HAVING} over such aggregates; \texttt{ORDER BY} names every
  grouping column, or every output \\
  Limits & \texttt{LIMIT} (at least one row) and \texttt{OFFSET} take integer literals; plain \texttt{DISTINCT}
  needs an \texttt{ORDER BY} over every output \\
  \bottomrule
  \end{tabular}
\end{table}


\paragraph{Scope attribution and resampling.}
Within each dataset, we order attempts randomly. We assign every oracle-admitted call the innermost scope in which an
exact-text equivalent occurred earlier. The scopes are the same session, another attempt of the same task and model, and another model on the same task. The outermost scopes are another task of the dataset and none. Shares are means over 200 random orders,
weighted by calls, by isolated DB-execution time and by payload bytes. To measure how reuse grows with the number of
attempts $N$, we sample $N$ attempts per task and model without replacement. We use 200 resamples for $N \in
\{1,2,4,8,16,50\}$. For each such batch we simulate a session cache (S1) and a cache shared by the batch's attempts,
both unbounded and starting empty. We label the latter S2. Without a budget, its DB-execution saving is the same for every cross-attempt design of Section~\ref{sec:designs}. Those designs differ in what they do with the result file. In all,
268 of the 270 task--model cells have the 50 runs that $N=50$ needs. We report medians over the 54 tasks with
task-bootstrap 95\% confidence intervals and check them with a bootstrap that resamples whole datasets. We test
paired differences with Wilcoxon signed-rank tests, Holm-corrected over a family of 29 tests. Hypotheses and success criteria were fixed in an experiment plan after a pilot replay of the SQLite and DuckDB calls. These are 64\% of the calls. The plan preceded the full replay, and Section~\ref{sec:prereg} reports the outcomes.

\paragraph{Instruments.}
A trace-driven simulator replays sampled attempt orders through each design of Section~\ref{sec:designs}. It accounts
the tool time, bytes written and peak footprint of every call. A live runner executes the same designs for real on
the calibration cells. These are 12 task--model cells chosen in advance to span each engine's inline-only, mixed and
spill-heavy tasks. The runner returns DAB's exact tool messages and writes real files. Its bytes written and peak
footprint match the simulator exactly in all 480 runs of the E-eval calibration. A live tool-boundary trace replay
runs the designs on one eight-attempt batch from each of the 270 task--model cells. It re-issues each attempt's
recorded database calls and re-runs its recorded Python consumers at their trace position against each design's
bindings. It also probes the shared store with write and delete attempts, and the LLM is not re-run. A threaded
stress test runs attempts concurrently against one shared store. A clean-room job re-parsed the original archive and
re-ran the scope and saving analyses. It reproduced all 8,546 values in their result files to a relative tolerance of
$10^{-9}$. Four further seeds leave the scope shares unchanged and move the $N=4$ median saving by at most 0.6 points.

\paragraph{Environment.}
Every replay, simulation and live run ran as a batch job in Singularity containers on CPU nodes of a university
cluster. The nodes are two-socket x86-64 machines with 28 or 36 cores and 180\,GiB of memory. Each job received the
cores and memory it requested, without exclusive use of its node. A replay executes each query up to three times in
a row after its dataset is loaded. We report the median of the completed executions, and Section~\ref{sec:robust}
weights by the first one instead. Jobs that share a node can slow a run, and the median of three executions damps
this. The live calibration runs every design of a cell within one job, in a random order per resample.
